% =====================================================================
% algebraic_islands.tex -- DRAFT v10 REVISION
%
% Prompted by Frembs, arXiv:2609.16490 (Leibniz Hannover), which cites this
% paper as [17] and uses it substantively in its Appendix A.
%
% SCOPE OF THE CORRECTION
% -----------------------
% No computational result in the paper changes.  Orthogonality was always
% decided by exact evaluation of the Hermitian product, never by reading the
% enumeration in Table~\ref{tab:two-element}; every ray count, pair count,
% triad count, island and OCUS-certified minimum stands.  What changes is the
% CLASSIFICATION NARRATIVE: the enumeration covers three-term vanishing sums
% only, while Definition~\ref{def:cancellation} admits sums of "at most three
% terms".  Two-term sums were never enumerated, and one of them -- Tr(x) = 0 --
% turns out to carry half the orthogonality structure of the Z[sqrt(-2)] pool.
%
% Verification for every number quoted below: ks_verify_two_term.py
% (exact arithmetic in Q(sqrt(d)), no floating point).  Positive control: the
% same script reproduces the published rows for sqrt(-2) (49/120/16),
% sqrt(-3) and sqrt(-5) (49/114/10), and the enriched Gaussian pool (127 rays,
% 51 triads) exactly.
%
% THREE ITEMS, in descending order of importance.
% =====================================================================


% ---------------------------------------------------------------------
% ITEM 1.  REPLACE Table~\ref{tab:two-element} (currently ~line 552).
% The table gains a two-term block; the three-term block is unchanged.
% ---------------------------------------------------------------------

\begin{table}[h]
\centering
\caption{All primitive vanishing sums from the product set $\bar{\mathcal{A}}
\cdot \mathcal{A}$ of a two-element alphabet $\{0, \pm 1, \pm x\}$, listed by
length.  A two-term sum suffices when one coordinate of the Hermitian product
vanishes; otherwise three terms are needed.  Each row is an algebraic constraint
on~$x$.  This classification applies to the raw alphabet before cross-product
completion.}
\label{tab:two-element}
\footnotesize
\setlength{\tabcolsep}{3pt}
\begin{tabular}{@{}llll@{}}
\toprule
Length & Zero-sum pattern & Identity on $x$ & Island / outcome \\
\midrule
2 & $x + \bar{x} = 0$        & $\mathrm{Tr}(x) = 0$   & structural$^\ddagger$ \\
2 & $1 - |x|^2 = 0$          & $|x|^2 = 1$            & roots of unity \\
2 & $1 - 1$, $x - x$, $|x|^2 - |x|^2$ & (none)        & trivial$^\S$ \\
\midrule
3 & $1 + 1 - x = 0$          & $x = 2$                & Integer \\
3 & $1 + 1 - |x|^2 = 0$      & $|x|^2 = 2$            & Peres/$\Z[\!\sqrt{-2}]$/Heeg.-7 \\
3 & $1 + x - |x|^2 = 0$      & $|x|^2 = 1 + x$        & Golden \\
3 & $1 + x + \bar{x} = 0$    & $\mathrm{Tr}(x) = -1$  & Eisenstein \\
3 & $1 - x - \bar{x} = 0$    & $\mathrm{Tr}(x) = +1$  & (colorable)$^\dagger$ \\
3 & $2|x|^2 - 1 = 0$         & $|x|^2 = \tfrac{1}{2}$ & $\equiv$ Peres (rescaled) \\
3 & $2x - |x|^2 = 0$         & $|x|^2 = 2x$           & $\equiv$ Integer \\
\bottomrule
\end{tabular}

\smallskip
\noindent $^\ddagger$\,Trace-zero cancellation is neither necessary nor sufficient
for KS-uncolorability, so it creates no island of its own: $\Z[i]$
($\mathrm{Tr}=0$, $|x|^2=1$) and $\Z[\sqrt{-3}]$, $\Z[\sqrt{-5}]$
($\mathrm{Tr}=0$, $|x|^2=3,5$) are colorable, while $\Z[\sqrt{2}]$
($\mathrm{Tr} = 2\sqrt2 \neq 0$) is not.  It nevertheless governs a large part of
the internal structure of the complex pools; see
Observation~\ref{obs:conjugation-dependence}.

\noindent $^\S$\,Available to every alphabet and independent of~$x$.  The
minimal alphabet $\{0,\pm1\}$ has only these, and generates 13 rays.

\noindent $^\dagger$\,The zero-sum $1 - x - \bar{x} = 0$ holds for all imaginary
quadratic generators with $d \equiv 1 \pmod{4}$.  It produces all-nonzero
orthogonal pairs and triads, but these are structurally decoupled from the
one-zero triads and the full ray set remains colorable.  The absence of type-B
(mixed) triads is proved for all integer $N \geq 3$ by SMT
verification~\cite{deMouraBjorner2008}.
\end{table}


% ---------------------------------------------------------------------
% ITEM 1b.  REPLACE Observation~\ref{obs:two-element} (currently ~line 575)
% and ADD the new Observation that carries the measurement.
% ---------------------------------------------------------------------

\begin{observation}\label{obs:two-element}
For any two-element alphabet $\{0, \pm 1, \pm x\}$, every primitive vanishing sum
in the product set $\bar{\mathcal{A}} \cdot \mathcal{A}$---of two terms or of
three---is listed in Table~\ref{tab:two-element}.  Of these, the patterns that
lead to KS-uncolorability reduce to one of the six non-cubic islands (up to
scaling equivalence).  Two further patterns produce genuine orthogonalities
without producing an island: $\mathrm{Tr}(x) = 1$ yields only colorable ray sets,
and $\mathrm{Tr}(x) = 0$ modulates the internal structure of a pool without
determining its colorability.
\end{observation}

\begin{observation}[Conjugation-dependence of the modulus-2 realizations]
\label{obs:conjugation-dependence}
The Peres, $\Z[\sqrt{-2}]$ and Gaussian realizations of modulus-2 cancellation
have isomorphic minimal 33-ray graphs (Proposition~\ref{prop:isomorphism}), but
they do not reach those graphs by the same arithmetic.  Call an orthogonality
\emph{conjugation-dependent} if the Hermitian product $\sum_k \bar v_k w_k$
vanishes while the bilinear product $\sum_k v_k w_k$ does not; such a pair exists
only because $\bar x \neq x$, which for a two-letter alphabet means
$\mathrm{Tr}(x) = 0$.  Counting exactly:

\medskip
\begin{center}
\begin{tabular}{lccc}
\hline
Pool & Pairs & Conj.-dependent & $\mathrm{Tr}$(gen.) \\
\hline
$\Z[\sqrt{2}]$, $\{0,\pm1,\pm\sqrt2\}$   & 120 & 0 \phantom{00}(0\%)  & $2\sqrt2$ \\
$\Z[\sqrt{-2}]$, $\{0,\pm1,\pm\sqrt{-2}\}$ & 120 & 60 \phantom{0}(50\%) & $0$ \\
$\Z[i]$, $\{0,\pm1,\pm i,\pm(1{+}i)\}$   & 465 & 327 (70\%) & $2$ \\
\hline
\end{tabular}
\end{center}
\medskip

\noindent Exactly half the orthogonality of the $\Z[\sqrt{-2}]$ pool---60 of 120
pairs, comprising 54 two-term and all 6 three-term cancellations---has no
counterpart in the real Peres pool, where the count is zero.  In the Gaussian
case the conjugation-dependence is supplied by the unit~$i$
($\mathrm{Tr}(i) = 0$, $|i|^2 = 1$) rather than by the modulus-2 generator
$1+i$, whose trace is~2.  The three rings therefore realize isomorphic graphs by
disjoint arithmetic routes.
\end{observation}


% ---------------------------------------------------------------------
% ITEM 2.  QUALIFY the "cancellation identity, not the ambient field" claim.
% It appears in at least four places.  Suggested replacement sentence, to be
% adapted at each site:
%
%   (i)   Prop.~\ref{prop:isomorphism}, closing sentence (~line 497)
%   (ii)  Obs.~\ref{obs:gaussian} (~line 488)
%   (iii) Sec.~\ref{sec:discussion}, end of the mechanism discussion (~line 523)
%   (iv)  Limitations item "Classification scope" (~line 924)
%
% The claim is not false, but it is stronger than the evidence supports: the
% identity determines the graph, while the ring retains arithmetic the graph
% does not record.
% ---------------------------------------------------------------------

The cancellation identity, not the ambient field, determines which orthogonality
\emph{graph} is realized.  It does not follow that the rings are interchangeable.
Frembs~\cite{Frembs2026} observes that $\Z[\sqrt{-2}]$ additionally satisfies the
two-term identity $\mathrm{Tr}(x) = 0$, which is unavailable at the real point
$x = \sqrt2$ and is not implied by $|x|^2 = 2$, and that over $\Z[\sqrt{-2}]$ this
relation generates an entire orbit of contexts and supports a 43-ray, 31-context
arrangement whose associated linear constraint system carries an obstruction
modulo~3 that the Peres geometry does not possess.  Observation
\ref{obs:conjugation-dependence} quantifies the difference: half the
orthogonality of the $\Z[\sqrt{-2}]$ pool is conjugation-dependent, against none
of the Peres pool.  Ordinary KS-uncolorability is therefore a coarser invariant
than the arithmetic of the coordinate ring, and graph isomorphism between two
islands does not certify that they support the same constructions.


% ---------------------------------------------------------------------
% ITEM 3.  SUSPECT ROW -- please re-run before v10 ships.
%
% Table~\ref{tab:complex-quadratic} (~line 326) records for Z[i], d = 1,
% alphabet {0, +-1, +-i}:   25 rays, 48 pairs, 6 triads.
%
% An independent reimplementation gives 31 rays, 69 pairs, 7 triads.
% Argument: every nonzero letter of {0,+-1,+-i} is a unit of Z[i], so any
% scalar carrying one alphabet word to another is a ratio of units, hence a
% unit.  The unit group {+-1, +-i} acts freely on the 5^3 - 1 = 124 nonzero
% words, giving 124/4 = 31 rays exactly.  (For d = 2,3,5 the generator is a
% non-unit, extra identifications occur, and both counts agree at 49.)
%
% The row's VERDICT is unaffected -- the pool is colorable either way and is
% not an island -- so this is a wrong number in a negative row, not a
% structural error.  Note also that the d = 1 row of Table~\ref{tab:heegner}
% uses the enriched alphabet {0,+-1,+-i,+-(1+i)} and its 127 rays / 51 triads
% reproduce exactly; only tab:complex-quadratic is in question.
% ---------------------------------------------------------------------


% ---------------------------------------------------------------------
% ITEM 4.  NEW "Note added in fourth revision", after the third-revision note.
% ---------------------------------------------------------------------

\subsection*{Note added in fourth revision}

\begin{sloppypar}
Frembs~\cite{Frembs2026}, in the course of classifying which inconsistent linear
constraint systems admit finite-dimensional quantum solutions, points out that the
enumeration of Table~\ref{tab:two-element} was confined to three-term vanishing
sums, whereas Definition~\ref{def:cancellation} admits sums of at most three
terms.  The omitted class is the two-term sums, and the productive member of that
class is $x + \bar x = 0$, the condition $\mathrm{Tr}(x) = 0$.  The table and
Observation~\ref{obs:two-element} have been extended accordingly.

No computational result is affected.  Orthogonality was decided throughout by
exact evaluation of the Hermitian product rather than by consulting the
enumeration, so the trace-zero cancellations were always counted; they were simply
not named.  Every ray count, pair count, triad count, island and OCUS-certified
minimum stands.  What the omission concealed is how much of the structure they
carry: exactly 60 of the 120 orthogonal pairs in the $\Z[\sqrt{-2}]$ pool vanish
only under conjugation, against 0 of the 120 in the real Peres pool
(Observation~\ref{obs:conjugation-dependence}).  The two pools have identical ray,
pair and triad counts and isomorphic graphs, but they arrive at them through
disjoint sets of identities.

This tempers, without overturning, the claim that the cancellation identity rather
than the ambient field controls the classification.  The claim holds for the
orthogonality graph, which is what the claim was tested against.  It does not
extend to everything a ring can support: Frembs exhibits a 43-ray, 31-context
arrangement over $\Z[\sqrt{-2}]$ built from the trace-zero relation, whose
associated constraint system is unsatisfiable over $\mathbb{F}_3$, while the
corresponding Peres incidence system admits no such nontrivial relation.  Ordinary
KS-uncolorability does not by itself produce that stronger obstruction.  We thank
Markus Frembs for the observation.
\end{sloppypar}


% ---------------------------------------------------------------------
% ITEM 5.  Bibliography, acknowledgments, reproducibility.
% ---------------------------------------------------------------------

\bibitem{Frembs2026}
M.~Frembs,
``A complete classification of the existence of finite-dimensional quantum
solutions to inconsistent linear constraint systems,''
arXiv:2609.16490 (2026).

% Acknowledgments -- add after the Cabello clause:
%   ``Markus Frembs for identifying the two-term omission in
%    Table~\ref{tab:two-element};''

% Reproducibility -- add to the script-to-table list:
%   \item \texttt{ks\_verify\_two\_term.py} $\to$
%         Obs.~\ref{obs:conjugation-dependence}, two-term rows of
%         Table~\ref{tab:two-element}

% ---------------------------------------------------------------------
% ITEM 6.  Unrelated bibliography fix noticed while preparing this draft.
%
% \bibitem{KirchwegarPeitlSzeider2023} reads ``A.~Peitl''.  The AAAI-23 author
% is Tom\'a\v{s} Peitl, so the initial should be ``T.''.  The citation key also
% misspells Kirchweger as ``Kirchwegar'' (harmless, but worth renaming).
% ---------------------------------------------------------------------
